\begin{explanationbox}
	\textbf{\textcolor{CExplanation}{一句话直觉}}
	圆锥内切球半径可由轴截面等腰三角形内切圆半径直接得出，利用等面积法 $R = \frac{S}{s}$ 一步到位，避免记忆复杂公式。

	\medskip
	\textbf{\textcolor{CExplanation}{核心拆解}}
	\begin{description}[style=nextline, leftmargin=2.8em, labelsep=0.5em, itemsep=0.45em]
		\item[\textbf{要点一（轴截面转化）：}]
			沿圆锥顶点和底面圆心作截面，得到等腰三角形，其内切圆半径即为球半径。
		\item[\textbf{要点二（等面积法）：}]
			等腰三角形面积 $S = \frac12\times 2r \times h = rh$，半周长 $s = r+l$，则 $R = S/s = \dfrac{rh}{r+l}$。
	\end{description}

	\medskip
	\textbf{\textcolor{CExplanation}{几何本质（截面法）}}
	内切球在轴截面上的大圆与三角形的三边相切，故球半径等于三角形内切圆半径，实现了三维到二维的转化。

	\medskip
	\textbf{\textcolor{CExplanation}{代数意义（数值关系）}}
	公式 $R = \dfrac{rh}{r+l}$ 中所有量均为正，且 $l = \sqrt{r^2+h^2}$，体现勾股关系。有理化变形 $R = \dfrac{r(l-r)}{h}$ 消去分母中的根式，方便试算。

	\medskip
	\textbf{\textcolor{CExplanation}{考点价值（高考方向）}}
	\begin{itemize}[leftmargin=*, itemsep=0.5em, topsep=0.4em]
		\item \textbf{考法一（直接套公式）：} 给出 $r,h$ 或 $r,l$ 直接求 $R$，只需代入公式。
		\item \textbf{考法二（逆向求尺寸）：} 已知 $R$ 及部分尺寸，通过公式反求 $r$ 或 $h$，常与勾股方程结合。
	\end{itemize}

	\medskip
	\textbf{\textcolor{CExplanation}{顿悟点}}
	核心在于“一切转化到轴截面”，把三维球问题变成二维三角形内切圆，等面积法是最快捷的路径。

	\medskip
	\textbf{\textcolor{CExplanation}{使用场景}}
	\begin{itemize}[leftmargin=*, itemsep=0.5em, topsep=0.4em]
		\item \textbf{场景一（直圆锥内切）：} 已知圆锥底面半径和高，直接求内切球半径。
		\item \textbf{场景二（组合体问题）：} 在球内切于圆锥的组合体中，利用该公式建立方程。
	\end{itemize}

\end{explanationbox}
